\part{ベクトル解析}
	\newcommand{\hth}{h_{\theta}}
	\newcommand{\hph}{h_{\varphi}}
	\newcommand{\dx}{\textrm{d}x}
	\newcommand{\dy}{\textrm{d}y}
	\newcommand{\dz}{\textrm{d}z}
	\newcommand{\ix}{\bm{i}_x}
	\newcommand{\iy}{\bm{i}_y}
	\newcommand{\iz}{\bm{i}_z}
	\newcommand{\ir}{\bm{i}_r}
	\newcommand{\ith}{\bm{i}_{\theta}}
	\newcommand{\iph}{\bm{i}_\varphi}
	\newcommand{\vr}{\bm{r}}
	\newcommand{\dc}{\textrm{d}\bm{c}}
	\newcommand{\dS}{\textrm{d}\bm{S}}
	\newcommand{\sinth}{\sin{\theta}}
	\newcommand{\costh}{\cos{\theta}}
	\newcommand{\sinph}{\sin{\varphi}}
	\newcommand{\cosph}{\cos{\varphi}}

	\newcommand {\dthv} [3] {\begin{bmatrix}#1\\#2\\#3\end{bmatrix}}	%3次元ベクトル

	\newcommand{\ppx}{\partDeriv{}{x}{}}
	\newcommand{\ppy}{\partDeriv{}{y}{}}
	\newcommand{\ppz}{\partDeriv{}{z}{}}
	\newcommand{\ppr}{\partDeriv{}{r}{}}
	\newcommand{\ppth}{\partDeriv{}{\theta}{}}
	\newcommand{\ppph}{\partDeriv{}{\varphi}{}}

	\chapter{3次元Euclid空間}
		\section{諸定義}
			\subsection{Vector Laplacian}
				\label{VectorLaplacian}
				\begin{definition}
					$A: \bm{r}\in\realNumbers^3 \mapsto A(\bm{r}) \in \realNumbers^3$($A$の各成分は$\mathrm{C}^2$級)の Vector Laplacian $\nabla^2 A$は次式で定義される。
					\[ \nabla^2A \coloneqq \nabla(\nabla\cdot A) - \nabla\times(\nabla\times A) \]
				\end{definition}
		\section{諸公式}
			\subsection{ベクトル三重積の公式: \texorpdfstring{$A\times(B\times C) = \inprod{A}{C}B - \inprod{A}{B}C$}{A×(B×C) = (A・C)B - (A・B)C}}
				\begin{proof}
					\quad\par
					$A,B,C$を次のようにおく。
					\begin{align*}
						A &= \bm{i}_1 A_1 + \bm{i}_2 A_2 + \bm{i}_3 A_3 \\
						B &= \bm{i}_1 B_1 + \bm{i}_2 B_2 + \bm{i}_3 B_3 \\
						C &= \bm{i}_1 C_1 + \bm{i}_2 C_2 + \bm{i}_3 C_3
					\end{align*}
					すると
					\begin{align*}
						A\times(B\times C) &= (\bm{i}_1 A_1 + \bm{i}_2 A_2 + \bm{i}_3 A_3) \times \left[\bm{i}_1(B_2 C_3 - B_3 C_2) + \bm{i}_2(B_3 C_1 - B_1 C_3) + \bm{i}_3(B_1 C_2 - B_2 C_1)\right] \\
						&= \phantom{+}\bm{i}_3 A_1(B_3 C_1 - B_1 C_3) - \bm{i}_2 A_1(B_1 C_2 - B_2 C_1) \\
						&\phantom{=} -\bm{i}_3 A_2(B_2 C_3 - B_3 C_2) + \bm{i}_1 A_2(B_1 C_2 - B_2 C_1) \\
						&\phantom{=} +\bm{i}_2 A_3(B_2 C_3 - B_3 C_2) - \bm{i}_1 A_3(B_3 C_1 - B_1 C_3) \\
						&= \phantom{+}\bm{i}_1 \left[(A_2 C_2 + A_3 C_3)B_1 - (A_2 B_2 + A_3 B_3)C_1\right] \\
						&\phantom{=} + \bm{i}_2 \left[(A_1 C_1 + A_3 C_3)B_2 - (A_1 B_1 + A_3 B_3)C_2\right] \\
						&\phantom{=} + \bm{i}_3 \left[(A_1 C_1 + A_2 C_2)B_2 - (A_1 B_1 + A_2 B_2)C_3\right] \\
						&= \phantom{+}\bm{i}_1 \left[(\inprod{A}{C} - A_1 C_1)B_1 - (\inprod{A}{B} - A_1 B_1)C_1\right] \\
						&\phantom{=} + \bm{i}_2 \left[(\inprod{A}{C} - A_2 C_2)B_2 - (\inprod{A}{B} - A_2 B_2)C_2\right] \\
						&\phantom{=} + \bm{i}_3 \left[(\inprod{A}{C} - A_3 C_3)B_3 - (\inprod{A}{B} - A_3 B_3)C_3\right] \\
						&= \phantom{+}\bm{i}_1 \left[\inprod{A}{C}B_1 - \inprod{A}{B}C_1\right] \\
						&\phantom{=} + \bm{i}_2 \left[\inprod{A}{C}B_2 - \inprod{A}{B}C_2\right] \\
						&\phantom{=} + \bm{i}_3 \left[\inprod{A}{C}B_3 - \inprod{A}{B}C_3\right] \\
						&= \inprod{A}{C}B - \inprod{A}{B}C
					\end{align*}
				\end{proof}
			\subsection{発散の別表現: \texorpdfstring{$\nabla\cdot A = \frac{1}{|V|}\lim_{|V|\to 0} \int_{\partial V} A(\bm{r})\cdot\bm{n}(\bm{r})\mathrm{d}^2\bm{r}$}{発散はベクトル場の法線面積分の体積の商の極限である}}
				\begin{shadebox}
					$\bm{r}_0\in\realNumbers^3$とし、$\Omega\subset\realNumbers^3$上のベクトル場$A$が$\bm{r}_0$の近傍で$C^1$級であるとする。
					$\bm{r}_0$を含む微小な領域を$V$とし、$\partial V$を$V$の境界、$|V|$を$V$の体積とする。
					但し$\partial V$は$C^1$級であるとする。
					このとき次式が成り立つ。
					\[ (\nabla\cdot A)(\bm{r}_0) = \frac{1}{|V|}\lim_{|V|\to 0}\int_{\partial V} A(\bm{r})\cdot\bm{n}(\bm{r})\mathrm{d}^2\bm{r} \]
					ここに$\bm{n}(\bm{r})$は$\bm{r}\in\partial V$に於ける外向きの法線ベクトルである。
				\end{shadebox}
				\begin{proof}
					\quad\par
					発散定理より次式が成り立つ。
					\[ \int_{\partial V} A(\bm{r})\cdot\bm{n}(\bm{r})\mathrm{d}^2\bm{r} = \int_V (\nabla\cdot A)(\bm{r})\mathrm{d}^3\bm{r} \]
					$\varepsilon>0$を任意に小さくとる。
					$A$が$\bm{r}_0$の近傍で$C^1$級であるという仮定から、$V$を充分小さくとれば、$|(\nabla\cdot A)(\bm{r}) - (\nabla\cdot A)(\bm{r}_0)| < \varepsilon$である。
					よって次式が成り立つ。
					\[ \left| \frac{1}{|V|}\int_V (\nabla\cdot A)(\bm{r})\mathrm{d}^3\bm{r} - (\nabla\cdot A)(\bm{r}_0) \right| < \varepsilon \]
					$|V|\to 0$とすれば定理の主張が成り立つ。
				\end{proof}
			\subsection{\texorpdfstring{$\nabla_\bm{r}\times f(\bm{r})\bm{A}(\bm{r}) = \bigl(\nabla_\bm{r}f(\bm{r})\bigr)\times \bm{A}(\bm{r}) + f(\bm{r})\nabla_\bm{r}\times\bm{A}(\bm{r})$}{ベクトルのスカラー関数倍の回転}}
				\label{ベクトルのスカラー関数倍の回転}
				\begin{shadebox}
					$\bm{r}\in\realNumbers$とする。
					$f: \realNumbers^3\to\complexNumbers,\;A(\bm{r}): \realNumbers^3\to\realNumbers^3$はともに$\mathrm{C}^1$級とする。
					このとき次式が成り立つ。
					\[ \nabla_\bm{r}\times f(\bm{r})\bm{A}(\bm{r}) = \bigl(\nabla_\bm{r}f(\bm{r})\bigr)\times \bm{A}(\bm{r}) + f(\bm{r})\nabla_\bm{r}\times\bm{A}(\bm{r}) \]
				\end{shadebox}
				\begin{proof}
					\quad\par
					$i+j$を3で割った余りを$[i+j]$と表すことにする。
					\begin{align*}
						\phantom{=} &\nabla_\bm{r}\times f(\bm{r})\bm{A}(\bm{r}) = \sum_{i=1}^3 \bm{i}_i\left(\partDeriv{f(\bm{r})A_{[i+2]}(\bm{r})}{r_{[i+1]}}{} - \partDeriv{f(\bm{r})A_{[i+1]}(\bm{r})}{r_{[i+2]}}{}\right) \\
						= &\sum_{i=1}^3 \bm{i}_i \left[\left(\partDeriv{f(\bm{r})}{r_{[i+1]}}{}A_{[i+2]}(\bm{r}) - \partDeriv{f(\bm{r})}{r_{[i+2]}}{}A_{[i+3]}(\bm{r})\right) + f(\bm{r})\left(\partDeriv{A_{[i+2]}(\bm{r})}{r_{[i+1]}}{} - \partDeriv{A_{[i+1]}(\bm{r})}{r_{[i+2]}}{}\right)\right] \\
						= &\bigl(\nabla_\bm{r}f(\bm{r})\bigr)\times \bm{A}(\bm{r}) + f(\bm{r})\nabla_\bm{r}\times\bm{A}(\bm{r})
					\end{align*}
				\end{proof}
			\subsection{3次元Euclid空間に於ける Vector Laplacian: \texorpdfstring{$\nabla^2 A = \bm{i}_1 \Delta A_1 + \bm{i}_2 \Delta A_2 + \bm{i}_3 \Delta A_3$}{∇\string^2 A = i1ΔA1 + i2ΔA2 + i3ΔA3}}
				\begin{shadebox}
					\ref{VectorLaplacian}に於いて、特に3次元Euclid空間を考える場合は次式のように簡約化される。
					\[ \nabla^2 A = \bm{i}_1 \Delta A_1 + \bm{i}_2 \Delta A_2 + \bm{i}_3 \Delta A_3 \]
				\end{shadebox}
				\begin{proof}
					\begin{align*}
						\nabla(\nabla\cdot A) &= \bm{i}_1 \partDerivLong{(\nabla\cdot A)}{r_1}{} + \bm{i}_2 \partDerivLong{(\nabla\cdot A)}{r_2}{} + \bm{i}_3 \partDerivLong{(\nabla\cdot A)}{r_3}{} \\
						&= \phantom{+} \bm{i}_1\left(\partDerivLong{A_1}{r_1}{2} + \partDerivIIHeteroLong{A_2}{r_1}{r_2} + \partDerivIIHeteroLong{A_3}{r_1}{r_3}\right) \\
						&\phantom{=} + \bm{i}_2\left(\partDerivIIHeteroLong{A_1}{r_2}{r_1} + \partDerivLong{A_2}{r_2}{2} + \partDerivIIHeteroLong{A_3}{r_2}{r_3}\right) \\
						&\phantom{=} + \bm{i}_3\left(\partDerivIIHeteroLong{A_1}{r_3}{r_1} + \partDerivIIHeteroLong{A_2}{r_3}{r_2} + \partDerivLong{A_3}{r_3}{2}\right)
					\end{align*}
					\begin{align*}
						\nabla\times(\nabla\times A) &= \nabla\times\left[\bm{i}_1\left(\partDeriv{A_3}{r_2}{} - \partDeriv{A_2}{r_3}{}\right) + \bm{i}_2\left(\partDeriv{A_1}{r_3}{} - \partDeriv{A_3}{r_1}{}\right) + \bm{i}_3\left(\partDeriv{A_2}{r_1}{} - \partDeriv{A_1}{r_2}{}\right)\right] \\
						&=\phantom{+} \bm{i}_1\left(-\partDerivLong{A_1}{r_2}{2} - \partDerivLong{A_1}{r_3}{2} + \partDerivIIHeteroLong{A_2}{r_2}{r_1} + \partDerivIIHeteroLong{A_3}{r_3}{r_2}\right) \\
						&\phantom{=} + \bm{i}_2\left(-\partDerivLong{A_2}{r_3}{2} - \partDerivLong{A_2}{r_1}{2} + \partDerivIIHeteroLong{A_3}{r_3}{r_2} + \partDerivIIHeteroLong{A_1}{r_1}{r_2}\right) \\
						&\phantom{=} + \bm{i}_3\left(-\partDerivLong{A_3}{r_1}{2} - \partDerivLong{A_3}{r_2}{2} + \partDerivIIHeteroLong{A_1}{r_1}{r_3} + \partDerivIIHeteroLong{A_2}{r_2}{r_3}\right)
					\end{align*}
					$A_1,A_2,A_3$が$\mathrm{C}^2$級であるから偏微分の順序が交換可能であるので、
					\[ \nabla(\nabla\cdot A) - \nabla\times(\nabla\times A) = \bm{i}_1 \Delta A_1 + \bm{i}_2 \Delta A_2 + \bm{i}_3 \Delta A_3 \]
				\end{proof}
			\subsection{\texorpdfstring{$\nabla\left(\inprod{\nabla f}{\bm{v}}\right) = \left(\nabla^2f\right)\bm{v}$}{∇(∇f v) = (∇\string^2 f)v}}
				\begin{shadebox}
					$\bm{v}$を定数ベクトルとするとき
					\[\nabla\left(\inprod{\nabla f}{\bm{v}}\right) = \left(\nabla^2f\right)\bm{v}\]
				\end{shadebox}
				\begin{proof}
					\quad\par
					\[\partDerivLong{\inprod{\nabla f}{\bm{v}}}{x_k}{} = \partDerivLong{\sum_{i=1}^{n}{\partDeriv{f}{x_i}v_i}}{x_k}{} = \sum_{i=1}^{n}{\partDerivIIHetero{f}{x_k}{x_i}}v_i\]
					であるから
					\[
						\nabla\left(\inprod{\nabla f}{\bm{v}}\right) = \left[
							\begin{array}{c}
								\partDerivLong{\inprod{\nabla f}{\bm{v}}}{x_1}{}\\
								\partDerivLong{\inprod{\nabla f}{\bm{v}}}{x_2}{}\\
								\vdots\\
								\partDerivLong{\inprod{\nabla f}{\bm{v}}}{x_n}{}
							\end{array}
						\right] = \left[
							\begin{array}{c}
								\left[\partDeriv{f}{x_1}{2}, \partDerivIIHetero{f}{x_1}{x_2},\dotsc,\partDerivIIHetero{f}{x_1}{x_n}\right]\bm{v}\\
								\vdots\\
								\left[\partDerivIIHetero{f}{x_n}{x_1}, \partDerivIIHetero{f}{x_n}{x_2},\dotsc,\partDeriv{f}{x_n}{2}\right]\bm{v}
							\end{array}
						\right] = \left(\nabla^2f\right)\bm{v}
					\]
				\end{proof}
			\subsection{\texorpdfstring{$\inprod{\bm{a}\times\bm{b}}{\bm{c}\times\bm{d}} = \inprod{\bm{a}}{\bm{c}}\inprod{\bm{b}}{\bm{d}} - \inprod{\bm{a}}{\bm{d}}\inprod{\bm{b}}{\bm{c}}$}{(a×b)・(c×d) = (a・c)(b・d) - (a・d)(b・c)}}
			\begin{proof}
				\quad\par
				ベクトル三重積の公式を利用する。
				\[ \inprod{\bm{a}\times\bm{b}}{\bm{c}\times\bm{d}} = \inprod{\bm{a}}{\bm{b}\times(\bm{c}\times\bm{d})} = \inprod{a}{\inprod{\bm{d}}{\bm{b}}\bm{c} - \inprod{\bm{b}}{\bm{c}}\bm{d}} = \inprod{\bm{a}}{\bm{c}}\inprod{\bm{b}}{\bm{d}} - \inprod{\bm{a}}{\bm{d}}\inprod{\bm{b}}{\bm{c}} \]
			\end{proof}
			\subsection{(系) \texorpdfstring{$\inprod{\bm{a}\times\bm{b}}{\bm{a}\times\bm{c}} = \norm{\bm{a}}^2\inprod{\bm{b}}{\bm{c}} - \inprod{\bm{a}}{\bm{b}}\inprod{\bm{a}}{\bm{c}}$}{(a×b)・(a×c) = ||a||\string^2(b・c) - (a・b)(a・c)}}
				%実は直前の公式よりもこちらの方を先に思いついた。直前の公式から直ちに導かれるが、当時思いついた証明を書き残しておく。
				\begin{proof}
					直前の公式から直ちに導ける。
					% \begin{align*}
					% 	\inprod{\bm{a}\times\bm{b}}{\bm{a}\times\bm{c}} &= (a_2b_1 - a_1b_2) (a_2c_1 - a_1c_2) + (a_3b_1 - a_1b_3) (a_3c_1 - a_1c_3) + (a_3b_2 - a_2b_3) (a_3c_2 - a_2c_3) \\
					% 	&=  {a_1}^2 b_2 c_2 + {a_1}^2 b_3 c_3 - a_1 a_2 b_1 c_2 - a_1 a_2 b_2 c_1 - a_1 a_3 b_1 c_3 - a_1 a_3 b_3 c_1 + {a_2}^2 b_1 c_1 \\
					% 	&\quad + {a_2}^2 b_3 c_3 - a_2 a_3 b_2 c_3 - a_2 a_3 b_3 c_2 + {a_3}^2 b_1 c_1 + {a_3}^2 b_2 c_2 \\
					% 	&= ({a_1}^2 + {a_2}^2 + {a_3}^2)(b_1c_1 + b_2c_2 + b_3c_3) - {a_1}^2b_1c_1 - {a_2}^2b_2c_2 - {a_3}^2b_3c_3 \\
					% 	&\quad - a_1 a_2 b_1 c_2 - a_1 a_2 b_2 c_1 - a_1 a_3 b_1 c_3 - a_1 a_3 b_3 c_1 - a_2 a_3 b_2 c_3 - a_2 a_3 b_3 c_2 \\
					% 	&= ({a_1}^2 + {a_2}^2 + {a_3}^2)(b_1c_1 + b_2c_2 + b_3c_3) - (a_1b_1 + a_2b_2 + a_3b_3)(a_1c_1 + a_2c_2 + a_3c_3) \\
					% 	&= \norm{\bm{a}}^2\inprod{\bm{b}}{\bm{c}} - \inprod{\bm{a}}{\bm{b}}\inprod{\bm{a}}{\bm{c}}
					% \end{align*}
				\end{proof}
			\subsection{\texorpdfstring{$\nabla_\bm{r}\frac{1}{\norm{\bm{r}+\bm{a}}} = -\frac{\bm{r}+\bm{a}}{\norm{\bm{r}+\bm{a}}^3}$}{距離に反比例するポテンシャルの勾配}}
				\begin{shadebox}
					$\bm{r} \coloneqq [r_1,r_2,r_3]^\top,\bm{a} \coloneqq [a_1,a_2,a_3]^\top\in\realNumbers^3$とする。次式が成り立つ。
					\[ \nabla_\bm{r}\frac{1}{\norm{\bm{r}+\bm{a}}} = -\frac{\bm{r}+\bm{a}}{\norm{\bm{r}+\bm{a}}^3} \]
				\end{shadebox}
				\begin{proof}
					\quad\par
					$i=1,2,3$に対して次式が成り立つ。
					\[ \partDerivLong{\frac{1}{\norm{\bm{r}+\bm{a}}}}{r_i}{} = \partDerivLong{\bigl((r_1+a_1)^2 + (r_2+a_2)^2 + (r_3+a_3)^2\bigr)^{-1/2}}{r_i}{} = -\frac{1}{2} \frac{1}{\norm{\bm{r}+\bm{a}}^3} \cdot 2(r_i+a_i) = -\frac{r_i+a_i}{\norm{\bm{r}+\bm{a}}^3} \]
				\end{proof}
			\subsection{\texorpdfstring{$\nabla_\bm{r} f(\bm{r})A = -A\times\nabla_\bm{r}f(\bm{r})$}{関数倍の定数ベクトルの回転}}
				\begin{shadebox}
					$\bm{r} \coloneqq [r_1,r_2,r_3]^\top,A \coloneqq [A_1,A_2,A_3]^\top\in\realNumbers^3,\;f:\realNumbers^3\to\complexNumbers$とする。次式が成り立つ。
					\[ \nabla_\bm{r} f(\bm{r})A = -A\times\nabla_\bm{r}f(\bm{r}) \]
				\end{shadebox}
				\begin{proof}
					\quad\par
					左辺の第1成分は次式で求まる。
					\[ \bigl(\nabla_\bm{r} f(\bm{r})A\bigr)[1] = \partDerivLong{A_3f(\bm{r})}{r_2}{} - \partDerivLong{A_2f(\bm{r})}{r_3}{} = A_3\partDerivLong{f(\bm{r})}{r_2}{} - A_2\partDerivLong{f(\bm{r})}{r_3}{} = -\bigl(A\times\nabla_\bm{r}f(\bm{r})\bigr)[1] \]
					同様のことを他の成分についても示すことができる。
				\end{proof}
		\section{Helmholtzの定理}
			\subsection{補題: 全空間に渡る積分が存在するための十分条件}
				\label{全空間に渡る積分が存在するための十分条件}
				\begin{shadebox}
					$\bm{r}\in\realNumbers^3$とする。
					$f:\realNumbers^3\to\complexNumbers$は十分遠方に於いて、その絶対値がある$\alpha>0$に対して$o(\norm{\bm{r}}_2^{-3-\alpha})$であるとする。
					このとき$\realNumbers^3$上での$f$の積分は存在する。
					すなわち次式が成り立つ。
					\[ R>0,\; V(R) \coloneqq \set{\bm{r}\in\realNumbers^3}{\norm{\bm{r}}_2 \leq R} ,\;\lim_{R\to\infty} \abs{\int_{V(R)} f(\bm{r}) \mathrm{d}^3\bm{r}} < \infty \]
				\end{shadebox}
				\begin{proof}
					\quad\par
					$R>0$は十分大きいものとし、$I(R)$を次式で定義する。
					\[ I(R) \coloneqq \int_{V(R)} f(\bm{r}) \mathrm{d}^3\bm{r} \]
					球座標変換$\bm{r} = r[\sin\theta\cos\phi,\;\sin\theta\sin\phi,\;\cos\theta]^\top$を用いて次式を得る。
					\[ I(R) = \integrate{0}{R}{\integrate{0}{\pi}{\integrate{0}{2\pi}{f(\bm{r})r^2\sin\theta}{}{\phi}}{}{\theta}}{}{r} \]
					$0<R_0<R$なる$R_0$をとり、$I(R) = I(R_0) + I_2(R)$と分解する。
					ここに$I_2(R) \coloneqq I(R) - I(R_0)$である。
					\begin{align*}
						|I_2(R)| &\leq \integrate{R_0}{R}{\integrate{0}{\pi}{\integrate{0}{2\pi}{|f(\bm{r})|r^2\sin\theta}{}{\phi}}{}{\theta}}{}{r} = 4\pi\integrate{R_0}{R}{o(\norm{\bm{r}}_2^{-3-\alpha})r^2}{}{r} \\
						&= 4\pi\integrate{R_0}{R}{o(r^{-1-\alpha})}{}{r} < \infty \text{ as } R\to\infty
					\end{align*}
				\end{proof}
			\subsection{補題: 与えられた発散を有し、回転が\texorpdfstring{$\bm{0}$}{0}であるベクトル場の構成}
				\label{与えられた発散を有し、回転が無いベクトル場の構成}
				Diracのデルタ関数を使わずに証明したいが、未達成。今後の課題。
			\subsection{補題: 与えられた回転を有し、発散が0であるベクトル場の構成}
				\label{与えられた回転を有し、発散が無いベクトル場の構成}
				Diracのデルタ関数を使わずに証明したいが、未達成。今後の課題。
			\subsection{Helmholtzの定理: 任意の\texorpdfstring{$\mathrm{C}^1$}{C\string^1}級のベクトル場は回転が\texorpdfstring{$\bm{0}$}{0}である\texorpdfstring{$\mathrm{C}^1$}{C\string^1}級の場と発散が0である\texorpdfstring{$\mathrm{C}^1$}{C\string^1}級の場に分解できる}
				任意に与えられた、発散および回転が十分遠方で0に収束するような$\mathrm{C}^1$級のベクトル場$A$が与えられたとき、補題\ref{与えられた発散を有し、回転が無いベクトル場の構成}により$A$と同じ発散を持ち、かつ回転が$\bm{0}$であるようなベクトル場$B$を構成できる。
				このとき$A - B$は発散が0であるので、補題\ref{与えられた回転を有し、発散が無いベクトル場の構成}により、回転がこれに一致するようなベクトル場$C$を構成できる。
				これらを用いて$A = B + C$と分解すればよい。
		\section{諸定理}
			\subsection{球殻定理}
				\begin{shadebox}
					$R>0,\;V \coloneqq \set{\bm{r}\in\realNumbers^3}{\norm{\bm{r}}_2\leq R}$とする。
					ベクトル場$\bm{A}(\bm{r})$を次式で定義する。
					\[ \bm{A}(\bm{r}) \coloneqq \int_V \frac{\bm{r}-\bm{r}'}{\norm{\bm{r}-\bm{r}'}_2^3} \mathrm{d}^3\bm{r}' \]
					このとき次が成り立つ。
					\[
						\bm{A}(\bm{r}) = \frac{4}{3}\pi
						\begin{cases}
							\bm{r} & (\norm{\bm{r}}_2\leq R) \\
							\frac{R^3}{\norm{\bm{r}}_2^3}\bm{r} & (\norm{\bm{r}}_2 > R)
						\end{cases}
					\]
				\end{shadebox}
				\begin{proof}
					\quad\par
					必要なら系全体を回転させ、一般性を失わず$\bm{r} = [0,0,r]^\top\;(r>0)$とする。
					球座標変換を用いて積分を実行する。
					$\bm{r}' \coloneqq r'[\sin\theta\cos\phi,\;\sin\theta\sin\phi,\;\cos\theta]^\top$とすると
					\begin{align*}
						\bm{A}(\bm{r}) &= \integrate{0}{R}{\integrate{0}{\pi}{\integrate{0}{2\pi}{\frac{\bm{r}-\bm{r}'}{\norm{\bm{r}-\bm{r}'}_2^3}r'^2(\sin\theta)}{}{\phi}}{}{\theta}}{}{r'} \\
						&= 2\pi\bm{i}_3\integrate{0}{R}{\integrate{0}{\pi}{\frac{r-r'\cos\theta}{(r^2+r'^2-2rr'\cos\theta)^{3/2}}r'^2(\sin\theta)}{}{\theta}}{}{r'}
					\end{align*}
					ここに$\bm{i}_3$は第3軸方向の単位ベクトルである。
					\par
					$f(r) \coloneqq -(r^2+r'^2-2rr'\cos\theta)^{-1/2}r'^2\sin\theta$とすると次式が成り立つ。
					\[ \deriv{f}{r}{}(r) = \frac{(r-r'\cos\theta)r'^2\sin\theta}{(r^2+r'^2-2rr'\cos\theta)^{3/2}} \]
					上式右辺の$(r',\theta)=(0,0)$近傍での振る舞いを考える。
					$r'=r+\Delta r\;(\Delta r\neq 0)$とすると、$\sin\theta = O(\theta),\;\cos\theta = 1+O(\theta^2)$なので分母・分子ともに$r^3 O(\theta^3)$となる。
					よって積分領域に於いて$f$は$r$の$\mathrm{C}^1$級関数であるから積分と微分の順序交換が可能であり、次式が成り立つ。
					\begin{align*}
						\bm{A}(\bm{r}) &= 2\pi\bm{i}_3\derivLong{\integrate{0}{R}{\integrate{0}{\pi}{f(r)}{}{\theta}}{}{r'}}{r}{} \\
						&= -2\pi\bm{i}_3\derivLong{\integrate{0}{R}{\integrate{0}{\pi}{\frac{r'^2\sin\theta}{\sqrt{r^2+r'^2-2rr'\cos\theta}}}{}{\theta}}{}{r'}}{r}{}
					\end{align*}
					$g(\theta) \coloneqq \frac{r'}{r}\sqrt{r^2+r'^2-2rr'\cos\theta}$とすると次式が成り立つ。
					\[ \deriv{g}{\theta}{}(\theta) = \frac{r'^2\sin\theta}{\sqrt{r^2+r'^2-2rr'\cos\theta}} \]
					よって次式が成り立つ。
					\[ \bm{A}(\bm{r}) = -2\pi\bm{i}_3\derivLong{\integrate{0}{R}{\left[g(\theta)\right]_0^\pi}{}{r'}}{r}{} = -2\pi\bm{i}_3\derivLong{\integrate{0}{R}{\frac{r'}{r}\bigl(|r+r'|-|r-r'|\bigr)}{}{r'}}{r}{} \]
					$R<r$のときは次式が成り立つ。
					\[ \bm{A}(\bm{r}) = -2\pi\bm{i}_3\derivLong{\integrate{0}{R}{\frac{2r'^2}{r}}{}{r'}}{r}{} = \frac{4\pi R^3}{3r^2}\bm{i}_3 = \frac{4\pi R^3}{3}\frac{\bm{r}}{\norm{\bm{r}}_2^3} \]
					$R\geq r$のときは$r'$に関する積分領域を$[0,r]$と$[r,R]$に分けて考える。
					前者については$R<r$のときの結果で$R$を$r$で置き換えたものとなる。
					後者については$|r+r'|-|r-r'| = 2r$となり、次式が成り立つ。
					\[ \bm{A}(\bm{r}) = -2\pi\bm{i}_3\derivLong{\integrate{0}{R}{2r'}{}{r'}}{r}{} = \bm{0} \]
				\end{proof}
				これを検証したMathematicaノートブックを\texttt{parts/vectorCalculus/mathematica/\\ shell\_theorem.nb}に置いてある。
			\subsection{非負領域の共通部分の非負領域}
				\begin{shadebox}
					$\realNumbers^n$上で考える。$\bm{v}_1,\dotsc,\bm{v}_m\;(m\leq n)$は一次独立であるとし、$\bm{v}$の非負領域を
					\[P(\bm{v})\coloneqq\left\{\bm{w}|\inprod{\bm{v}}{\bm{w}}\geq0\right\})\]
					で定義する。あるベクトル$\bm{u}$が$\bigcap_{i=1}^{m}{P(\bm{v}_i)}$の非負領域にある、すなわち
					\[\inprod{\bm{u}}{\bm{w}}\geq0,\forall\bm{w}\in\bigcap_{i=1}^{m}{P(\bm{v}_i)}\]
					であるとき、$\bm{u}$は適当な定数$a_i\geq0(i=1,\dotsc,m)$によって
					\[\bm{u} = \sum_{i=1}^{m}{a_i\bm{v}_i}\]
					と表せる。
				\end{shadebox}
				\begin{proof}
					\quad\par
						まず$\bm{u}$が$\bm{v}_1,\dotsc,\bm{v}_m$の一次結合で表せることを示す。\\
						$W \coloneqq \Span{\bm{v}_1,\dotsc,\bm{v}_m}$とすると$\realNumbers^n = W\oplus W^\bot$であるから、適当な定数$c_1,\dotsc,c_m$及び適当なベクトル$\bm{w}_t\in W^\bot$を用いて
						\begin{equation}
							\bm{u} = \sum_{i=1}^{m}{c_i\bm{v}_i} + \bm{w}_t
						\end{equation}
						と表せる。$\bm{p}$を$\bigcap_{i=1}^{m}{P(\bm{v}_i)}$の任意のベクトルとすると
						\[\inprod{\bm{v}_i}{\bm{p}}\geq0,\forall\bm{v}_i\in W\]
						また、$\bm{w}_{t2}$を$W^\bot$の任意のベクトルとすると
						\[\inprod{\bm{v}_i}{\bm{w}_{t2}}=0,\forall\bm{v}_i\in W\]
						従って$\bm{p}' \coloneqq \bm{p}+\bm{w}_{t2}$とすると
						\[\inprod{\bm{v}_i}{\bm{p}'}\geq0,\forall\bm{v}_i\in W\]
						であるから$\bm{p'}$もまた$\bigcap_{i=1}^{m}{P(\bm{v}_i)}$のベクトルである。従って$\inprod{\bm{u}}{\bm{p}'}\geq0$であるから
						\[0\leq\inprod{\bm{u}}{\bm{p}'} = \inprod{\sum_{i=1}^{m}{c_i\bm{v}_i}}{\bm{p}} + \inprod{\bm{w}_{t1}}{\bm{w}_{t2}}\]
						である。然るに$\bm{w}_{t2}\in W^\bot$は任意であるから、$\bm{w}_{t1}=\bm{0}$でない限り$\inprod{\bm{w}_{t1}}{\bm{w}_{t2}}$の値をいくらでも小さくできる。
						よって上式が成り立つためには$\bm{w}_{t1}=\bm{0}$でなくてはならならず、式(1)より結局
						\[\bm{u} = \sum_{i=1}^{m}{c_i\bm{v}_i}\]
						と表せることになる。
						\par
						次に$c_i$が全て非負であることを示す。%そのために$\left(\bm{d}_i\in\bigcap_{i=1}^{m}{P(\bm{v}_i)}\right)\land\inprod{\bm{v}_j}{\bm{d}_i}=\delta_{ij}\;(\delta_{ij}\text{はクロネッカーのデルタ})$なるベクトル$\bm{d}_i$が存在することを示そう。
						$\bm{v}_1,\dotsc,\bm{v}_m\;(m\leq n)$は一次独立であったから次の行列
						\[
							A = \left[
								\begin{array}{c}
									{\bm{v}_1}^\textrm{T}\\
									{\bm{v}_2}^\textrm{T}\\
									\vdots\\
									{\bm{v}_m}^\textrm{T}
								\end{array}
							\right]
						\]
						の階数は$m$である。従って次の方程式
						\[
							A\bm{x} = \left[
								\begin{array}{c}
									0\\
									\vdots\\
									1\\
									\vdots\\
									0
								\end{array}
							\right]
							\begin{array}{c}
								\quad\\
								\quad\\
								\leftarrow i\\
								\quad\\
								\quad
							\end{array}
						\]
						は解を持つ。この解を$\bm{d}_i$とすれば$\inprod{\bm{v}_j}{\bm{d}_i}=\delta_{ij}$であるから$\bm{d}_i\in\bigcap_{i=1}^{m}{P(\bm{v}_i)}$。
						よって
						\[0\leq\inprod{\bm{u}}{\bm{d}_i} = c_i\]
				\end{proof}
		\subsection{Stokesの定理}
			\begin{shadebox}
				$c$を空間上の任意の$C^1$級の閉曲線、$S$を、$c$を周$\partial S$とする任意の$C^2$級の曲面, $A$を$C^1$級のベクトル場とするとき
				\[
					\oint_c{A\cdot\dc} = \iint_S{\nabla\times A}\cdot\dS
				\]
			\end{shadebox}
			\begin{proof}
				\quad\par
				$S$が$S(x,y)$の如く$x,y$の関数で表せる場合について証明すれば十分である。
				なぜなら、$S$が複雑な場合は$S(s,y),S(y,z),S(z,x)$いずれかの型の2変数関数で表せる曲面に分割して、各部分で、これから証明する$S(x,y)$型の場合と同じ手続きを踏んで最後に合計すれば全体の結果になるからである。
				領域分割を行う場合は分割境界でも線積分を行うことになるが、その回数は2で、互いに逆向きであるから、合計する段階で相殺して0になるため最終的な結果には現れない。
				\begin{figure}[H]
					\centering
					\includegraphics[keepaspectratio, scale=0.4]
					{\currfiledir/imgs/Stokes-theorem/divide.eps}
					\caption{$S$の分割}
				\end{figure}
				以降$S$は$S(x,y)$型の曲面とする。$\partial S$を$t$をパラメータとする閉曲線とする。
				すなわち
				\[\partial S\;:\;x=\phi(t), y=\psi(t), t\;:\;0\rightarrow1\]
				また、$xy$平面への$S$の射影を$\Omega$とする。
				\begin{figure}[H]
					\centering
					\includegraphics[keepaspectratio, scale=0.4]
					{\currfiledir/imgs/Stokes-theorem/illust.eps}
					\caption{状況}
				\end{figure}
				では証明しよう。
				線積分で表された定理の左辺を変形して面積分に置き換えて、定理の右辺と見比べる流れで行く。
				\par
				関数の引数を全部書いていてはスペースがいくらあっても足りないので省略して書くことになる(例えば$A_x(x,y,z)$を単に$A_x$と書いたり)のだが、こうすると独立変数と従属変数の関係を見誤る危険が生じるため注意が必要だ。
				迂闊に計算すると罠にはまる。
				特に危ない所では敢えて引数を明示する。
				\begin{align*}
					\oint_c{A\cdot\dc} &= \oint_{\partial S}{\left(\ix A_x+\iy A_y+\iz A_z\right)\cdot\left(\ix \dx+\iy \dy+\iz \dz\right)} \\
					&= \oint_{\partial S}{\left(A_x\dx+A_y\dy\right)} + \oint_{\partial S}{A_z\dz} \\
					&= \underbrace{\oint_{\partial \textcolor{red}{\Omega}}{\left(A_x(x,y,S(x,y))\dx+A_y(x,y,S(x,y))\dy\right)}}_{\ctext{A}} + \underbrace{\oint_{\partial S}{A_z\dz}}_{\ctext{B}}
				\end{align*}
				\ctext{A}で積分経路が$\partial S$から$\partial\Omega$に変更されたことに注意。Greenの定理を使いたいから$xy$平面に落としたのだ。
				Greenの定理より
				\[
					\ctext{A} = \iint_\Omega{\left[\underbrace{\partDerivLong{A_y(x,y,S(x,y))}{x}{}}_{\ctext{C}}-\underbrace{\partDerivLong{A_x(x,y,S(x,y))}{y}{}}_{\ctext{D}}\right]\dx\dy}
				\]
				\ctext{C}が罠だ。引数を明示していないと$\left.\partDeriv{A_y}{x}{}\right|_{(x,y,z)=(x,y,S(x,y))}$と勘違いしてしまう。
				\textcolor{red}{$\partDerivLong{A_y(x,y,S(x,y))}{x}{}$と$\left.\partDeriv{A_y}{x}{}\right|_{(x,y,z)=(x,y,S(x,y))}$は別物}であることに注意せなばならぬ。
				連鎖律に注意して
				\begin{align*}
					\ctext{C} &= \left.\partDeriv{A_y}{x}{}\right|_{(x,y,z)=(x,y,S(x,y))}\partDeriv{x}{x}{} + \left.\partDeriv{A_y}{y}{}\right|_{(x,y,z)=(x,y,S(x,y))}\partDeriv{y}{x}{} + \left.\partDeriv{A_y}{z}{}\right|_{(x,y,z)=(x,y,S(x,y))}\partDeriv{S}{x}{} \\
					&= \partDeriv{A_y}{x}{} + \partDeriv{A_y}{z}{}\partDeriv{S}{x}{}
				\end{align*}
				同様にして
				\[\ctext{D} = \partDeriv{A_x}{y}{} + \partDeriv{A_x}{z}{}\partDeriv{S}{y}{}\]
				よって結局
				\begin{align*}
					\ctext{A} &= \iint_\Omega{\left[\partDeriv{A_y}{x}{} + \partDeriv{A_y}{z}{}\partDeriv{S}{x}{} - \partDeriv{A_x}{y}{} - \partDeriv{A_x}{z}{}\partDeriv{S}{y}{}\right]\dx\dy} \\
					&= \iint_\Omega{\left[\left(\partDeriv{A_y}{x}{}-\partDeriv{A_x}{y}{}\right) + \partDeriv{A_y}{z}{}\partDeriv{S}{x}{} - \partDeriv{A_x}{z}{}\partDeriv{S}{y}{}\right]\dx\dy}
				\end{align*}
				次に\ctext{B}を見る。$t$による偏微分を$\dot{}$で表すことにする。線積分の定義に従って
				\begin{align*}
					\ctext{B} &= \integrate{0}{1}{A_z(\phi(t),\psi(t),S(\phi(t),\psi(t))\left(\partDeriv{S}{x}{}\dot{\phi(t)}+\partDeriv{S}{y}{}\dot{\psi(t)}\right)}{}{t} \\
					&= \integrate{0}{1}{\left(A_z(\phi(t),\psi(t),S(\phi(t),\psi(t))\partDeriv{S}{x}{}\right)\dot{\phi(t)}}{}{t} + \integrate{0}{1}{\left(A_z(\phi(t),\psi(t),S(\phi(t),\psi(t))\partDeriv{S}{y}{}\right)\dot{\psi(t)}}{}{t} \\
					&= \oint_{\partial \Omega}{A_z(x,y,S(x,y))\partDeriv{S}{x}{}\dx} + \oint_{\partial \Omega}{A_z(x,y,S(x,y))\partDeriv{S}{y}{}\dy} \\
					&= \oint_{\partial \Omega}{\left(A_z(x,y,S(x,y))\partDeriv{S}{x}{}\dx + A_z(x,y,S(x,y))\partDeriv{S}{y}{}\dy\right)}
				\end{align*}
				これにGreenの定理を適用して、連鎖律に注意して偏微分していく。
				\begin{align*}
					\ctext{B} &= \iint_\Omega{\left[\partDerivLong{\left(A_z(x,y,S(x,y))\partDeriv{S}{y}{}\right)}{x}{} - \partDerivLong{\left(A_z(x,y,S(x,y))\partDeriv{S}{x}{}\right)}{y}{}\right]\dx\dy} \\
					&= \iint_\Omega{\left(\partDeriv{A_z}{x}{}\partDeriv{S}{y}{} + \partDeriv{A_z}{z}{}\partDeriv{S}{x}{}\partDeriv{S}{y}{} + A_z\partDerivIIHetero{S}{y}{x} - \partDeriv{A_z}{y}{}\partDeriv{S}{x}{} - \partDeriv{nA_z}{z}{}\partDeriv{S}{y}{}\partDeriv{S}{x}{} - A_z\partDerivIIHetero{S}{x}{y}\right)\dx\dy} \\
					&= \iint_\Omega{\left(\partDeriv{A_z}{x}{}\partDeriv{S}{y}{} - \partDeriv{A_z}{y}{}\partDeriv{S}{x}{}\right)\dx\dy}\quad\left(\because\;\text{$S$は$C^2$級と仮定したから$\partDerivIIHetero{S}{y}{x} = \partDerivIIHetero{S}{x}{y}$}\right)
				\end{align*}
				以上より
				\begin{align*}
					\oint_c{A\cdot\dc} &= \iint_\Omega{\left[-\left(\partDeriv{A_z}{y}{}-\partDeriv{A_y}{z}{}\right)\partDeriv{S}{x}{}-\left(\partDeriv{A_x}{z}{}-\partDeriv{A_z}{x}{}\right)\partDeriv{S}{y}{}+\left(\partDeriv{A_y}{x}{}-\partDeriv{A_x}{y}{}\right)\right]\dx\dy}
				\end{align*}
				\par
				これが定理の右辺と等しいことを示せば良い。
				まず$\dS$の正体を暴こう。感覚的には、$\dS$とは、$S$上で$x,y$が微小量$\dx,\dy$だけ動いた時に舐める$S$の面積をノルムに持ち、向きはその微小面上の任意の代表点での上向きの法線ベクトルに一致するようなベクトルであった。
				所謂「面積ベクトル」である。
				\par
				厳密にはRiemann和を取る段階で定義され、$S$上の点$(x,y,S(x,y))$における接平面のうち、その射影が$xy$平面上の矩形$\dx\dy$に収まる部分(3次元上の平行四辺形)の面積をノルムに持ち、向きは接平面の上向き法線ベクトルと一致するようなベクトルである。
				\[\dS = \left(\ix\dx+\iz\partDeriv{S}{x}{}\dx\right)\times\left(\iy\dy+\iz\partDeriv{S}{y}{}\dy\right) = \left(-\ix\partDeriv{S}{x}{} - \iy\partDeriv{S}{y}{} + \iz\right)\dx\dy\]
				よって定理の右辺は
				\begin{align*}
					\iint_S{\nabla\times A}\cdot\dS &= \iint_\Omega{\left[\ix\left(\partDeriv{A_z}{y}{}-\partDeriv{A_y}{z}{}\right)+\iy\left(\partDeriv{A_x}{z}{}-\partDeriv{A_z}{x}{}\right)+\iz\left(\partDeriv{A_y}{x}{}-\partDeriv{A_x}{y}{}\right)\right]\cdot\dS} \\
					&= \iint_\Omega{\left[-\left(\partDeriv{A_z}{y}{}-\partDeriv{A_y}{z}{}\right)\partDeriv{S}{x}{}-\left(\partDeriv{A_x}{z}{}-\partDeriv{A_z}{x}{}\right)\partDeriv{S}{y}{}+\left(\partDeriv{A_y}{x}{}-\partDeriv{A_x}{y}{}\right)\right]\dx\dy} \\
					&= \oint_c{A\cdot\dc}
				\end{align*}
			\end{proof}
	\chapter{一般の直交座標系と3次元Euclid空間の関係}
		\section{座標変換}
			$T:\realNumbers^3\to\realNumbers^3$は可逆であり$\mathrm{C}^1$級であるとする。
			$T$のJacobi行列の列は互いに直交するものとする。
			$\bm{u}\in\realNumbers^3,\;\bm{r} = T(\bm{u})$とする。
			例えば球座標系では$\bm{u} = (r,\theta,\varphi)$であり、$\bm{r} = \bm{i}_x x + \bm{i}_y y + \bm{i}_z z,\;(x,y,z) = (r\sin\theta\cos\varphi,r\sin\theta\sin\varphi,r\cos\theta)$である。
			次式で定める量を$\bm{u}$に対する$\bm{r}$の計量係数(Lam\'{e} coefficient)と呼ぶ。
			\[ h_i \coloneqq \norm{\partDeriv{\bm{r}}{u_i}{}}_2 \quad (i=1,2,3) \]
			$h_i$は$T$のJacobi行列$J$の第$i$列のノルムであり、逆写像定理から$J$は可逆なので$h_i \neq 0$である。
			\par
			次式で定める単位ベクトルを「$u_i$方向単位ベクトル」と呼ぶ。
			\[ \bm{j}_i := \frac{1}{h_i}\partDeriv{\bm{r}}{u_i}{} \]
			$T$のJacobi行列の列が互いに直交するため、$\bm{j}_1,\bm{j}_2,\bm{j}_3$は互いに直交する。
		\section{表記の慣習上の注意}
			$\realNumbers^n$という記号は複数の文脈で使われる。
			ある時は単に$\realNumbers$の直積集合を表す記号として、またある時は$n$次元Euclid空間を表す記号として使われる。
			例えば球座標系の$(r,\theta,\varphi)$が属する構造は3次元Euclid空間ではなく単に$\realNumbers$の3つの直積集合であり、その$T$による像$(x,y,z)$もまた$\realNumbers$の3つの直積集合上の量である。
			そして$(x,y,z)$と3次元Euclid空間上のベクトル$\bm{i}_x x + \bm{i}_y y + \bm{i}_z z$は一対一で対応するから、しばしばこれらを同一視する。
		\section{単位ベクトルの変換}
			\label{単位ベクトルの変換}
			\begin{shadebox}
				既出の記号の定義は引き継ぐ。
				3次元Euclid空間の標準単位ベクトルは直交座標系の単位ベクトルを用いて次のように表せる。
				\[ \bm{i}_i = \sum_{j=1}^3 \frac{1}{h_j}\frac{\partial r_i}{\partial u_j}\bm{j}_j \quad (i=1,2,3) \]
				つまり行列を用いて次のように表せる。
				\[ [\bm{i}_1, \bm{i}_2, \bm{i}_3] = [\bm{j}_1, \bm{j}_2, \bm{j}_3]JH^{-1} \]
				ここに$J$は$T$のJacobi行列であり、$H := \diag{h_1,h_2,h_3}$である。
			\end{shadebox}
			\begin{proof}
				\quad\par
				$\bm{j}_1,\bm{j}_2,\bm{j}_3$は$\realNumbers^3$の正規直交基底であるから次式が成り立つ。
				\[ \bm{i}_i = \sum_{j=1}^3 (\bm{i}_i\cdot\bm{j}_j)\bm{j}_j = \sum_{j=1}^3 \left(\bm{i}_i\cdot\frac{1}{h_j}\frac{\partial \bm{r}}{\partial u_j}\right)\bm{j}_j = \sum_{j=1}^3 \left(\bm{i}_i\cdot\frac{1}{h_j}\sum_{k=1}^3\frac{\partial r_k}{\partial u_j}\bm{i}_k\right)\bm{j}_j \\
				= \sum_{j=1}^3 \frac{1}{h_j}\frac{\partial r_i}{\partial u_j}\bm{j}_j \]
			\end{proof}
		\section{ベクトルの成分の変換}
			\label{ベクトルの成分の変換}
			\begin{shadebox}
				既出の記号の定義は引き継ぐ。
				3次元Euclid空間上のベクトル$\bm{A}(\bm{r}) = \bm{i}_1 A_1(\bm{r}) + \bm{i}_2 A_2(\bm{r}) + \bm{i}_3 A_3(\bm{r})$を$u_1,u_2,u_3$直交座標系で表したものを$\bm{j}_1 \tilde{A}_1(\bm{u}) + \bm{j}_2 \tilde{A}_2(\bm{u}) + \bm{j}_3 \tilde{A}_3(\bm{u})$とするとき、$A_1,A_2,A_3$と$\tilde{A}_1,\tilde{A}_2,\tilde{A}_3$の間には次の関係式が成り立つ。
				\[ [A_1,A_2,A_3]^\top = JH^{-1}[\tilde{A}_1,\tilde{A}_2,\tilde{A}_3]^\top \]
				ここに$J$は$T$のJacobi行列であり、$H := \diag{h_1,h_2,h_3}$である。
			\end{shadebox}
			\begin{proof}
				\begin{align*}
					\phantom{=} &[\bm{i}_1,\bm{i}_2,\bm{i}_3][A_1,A_2,A_3] = \bm{A}(\bm{r}) = \bm{j}_1 \tilde{A}_1 + \bm{j}_2 \tilde{A}_2 + \bm{j}_3 \tilde{A}_3 = [\bm{j}_1,\bm{j}_2,\bm{j}_3][\tilde{A}_1,\tilde{A}_2,\tilde{A}_3]^\top \\
					= &[\bm{i}_1,\bm{i}_2,\bm{i}_3]JH^{-1}[\tilde{A}_1,\tilde{A}_2,\tilde{A}_3]^\top
				\end{align*}
			\end{proof}
		\section{勾配}
			\label{勾配}
			\begin{shadebox}
				既出の記号の定義は引き継ぐ。
				$\phi:\realNumbers^3\to\complexNumbers$は$\realNumbers^3$上で$\mathrm{C}^1$級であるとする。
				$\tilde{\phi}(\bm{u}) \coloneqq \phi(T(\bm{u}))$とする。
				$\phi$の勾配$\nabla_\bm{r}\phi(\bm{r})$は$\bm{u}$を用いて次式で表せる。
				\[ \nabla_\bm{r}\phi(\bm{r}) = \sum_{i=1}^3\frac{\bm{j}_i}{h_i}\partDeriv{\phi}{u_i}{} = \left[\frac{\bm{j}_1}{h_1}, \frac{\bm{j}_2}{h_2}, \frac{\bm{j}_3}{h_3}\right]^\top\nabla_\bm{u}\tilde{\phi}(\bm{u}) \]
			\end{shadebox}
			\begin{proof}
				\begin{gather*}
					\bm{j}_i = \frac{1}{h_i}\partDeriv{\bm{r}}{u_i}{} = \frac{1}{h_i}\partDerivLong{\sum_{j=1}^3 r_j\bm{i}_j}{u_i}{} = \frac{1}{h_i}[\bm{i}_1,\bm{i}_2,\bm{i}_3]\left[\partDeriv{r_1}{u_i}{},\partDeriv{r_2}{u_i}{},\partDeriv{r_3}{u_i}{}\right]^\top \\
					\therefore\;[\bm{j}_1,\bm{j}_2,\bm{j}_3] = [\bm{i}_1,\bm{i}_2,\bm{i}_3]\left[\frac{1}{h_j}\partDeriv{r_i}{u_j}{}\right]_{(i,j)\in\{1,2,3\}^2} = JH^{-1}
				\end{gather*}
				ここに$\quad H \coloneqq \diag{h_1,h_2,h_3}$であり、$J$は$T$のJacobi行列である。
				以上より次式が成り立つ。
				\[ [\bm{i}_1,\bm{i}_2,\bm{i}_3] = [\bm{j}_1,\bm{j}_2,\bm{j}_3]HJ^{-1} \]
				$\partDeriv{\phi}{r_i}{}$を評価すると次式を得る。
				\begin{gather*}
					\partDeriv{\phi}{r_i}{} = \sum_{j=1}^3\partDeriv{\tilde{\phi}}{u_j}{}\partDeriv{u_j}{r_i}{} \\
					\begin{split}
						\therefore \quad \nabla_\bm{r}\phi(\bm{r}) &= \left[\partDeriv{\phi}{r_1}{}, \partDeriv{\phi}{r_2}{}, \partDeriv{\phi}{r_3}{}\right]^\top = \left[\partDeriv{u_j}{r_i}{}\right]_{(i,j)\in\{1,2,3\}^2}\left[\partDeriv{\tilde{\phi}}{u_1}{},\partDeriv{\tilde{\phi}}{u_2}{},\partDeriv{\tilde{\phi}}{u_3}{}\right]^\top \\
						&= \left(J^{-1}\right)^\top\nabla_\bm{u}\tilde{\phi}(\bm{u})
					\end{split}
				\end{gather*}
				最後の等号は$\bm{u} = T^{-1}(\bm{r})$と逆関数定理による。
				以上より、次式が成り立つ。
				\[ \nabla_\bm{r}\phi(\bm{r}) = [\bm{j}_1,\bm{j}_2,\bm{j}_3]HJ^{-1}\left(J^{-1}\right)^\top\nabla_\bm{u}\tilde{\phi}(\bm{u}) \]
				$J^{-1}\left(J^{-1}\right)^\top = J^{-1}\left(J^\top\right)^{-1} = \left(J^\top J\right)^{-1}$であり、$T$に関する仮定から$J$の列は互いに直交し、計量係数の定義から$J^\top J$の第$i$対角成分は$h_i^2$である。
				よって$\left(J^\top J\right)^{-1} = H^{-2}$であり、次式が成り立つ。
				\[ \nabla_\bm{r}\phi(\bm{r}) = [\bm{j}_1,\bm{j}_2,\bm{j}_3]H^{-1}\nabla_\bm{u}\tilde{\phi}(\bm{u}) = \left[\frac{\bm{j}_1}{h_1}, \frac{\bm{j}_2}{h_2}, \frac{\bm{j}_3}{h_3}\right]^\top\nabla_\bm{u}\tilde{\phi}(\bm{u}) \]
			\end{proof}
		\section{発散}
			\label{発散}
			\begin{shadebox}
				既出の記号の定義は引き継ぐ。
				さらに$T$は$\textrm{C}^2$級であるとし、$\frac{\partial h_i}{\partial u_i} = 0\;(i=1,2,3)$が成り立つとする(例: 円柱座標系, 球座標系)。
				3次元Euclid空間上の$\mathrm{C}^1$級のベクトル場$\bm{A}(\bm{r}) = \bm{i}_1 A_1(\bm{r}) + \bm{i}_2 A_2(\bm{r}) + \bm{i}_3 A_3(\bm{r})$を$u_1,u_2,u_3$直交座標系で表したものを$\bm{j}_1 \tilde{A}_1(\bm{u}) + \bm{j}_2 \tilde{A}_2(\bm{u}) + \bm{j}_3 \tilde{A}_3(\bm{u})$とするとき、次式が成り立つ。
				\[ \nabla_\bm{r}\cdot\bm{A}(\bm{r}) = \frac{1}{h_1 h_2 h_3}\left[ \partDeriv{h_2 h_3 \tilde{A}_1(\bm{u})}{u_1}{} + \partDeriv{h_3 h_1 \tilde{A}_2(\bm{u})}{u_2}{} + \partDeriv{h_1 h_2 \tilde{A}_3(\bm{u})}{u_3}{} \right] \]
			\end{shadebox}
			\begin{proof}
				\[ \nabla_\bm{r}\cdot\bm{A}(\bm{r}) = \sum_{i=1}^3 \partDeriv{A_i(\bm{r})}{r_i}{} = \sum_{i=1}^3 \sum_{j=1}^3 \partDeriv{A_i(\bm{r})}{u_j}{}\partDeriv{u_j}{r_i}{} = \sum_{i=1}^3 \sum_{j=1}^3 \left(\partDerivLong{\sum_{k=1}^3\frac{1}{h_k}\partDeriv{r_i}{u_k}{}\tilde{A}_k(\bm{u})}{u_j}{}\right)\partDeriv{u_j}{r_i}{} \tag{1} \]
				最後の等号成立は\ref{ベクトルの成分の変換}による。
				\par
				まず$\partDeriv{u_j}{r_i}{}$を評価する。
				$J$を$T$のJacobi行列とすると次式が成り立つ。
				\[ \partDeriv{u_j}{r_i}{} = J^{-1}[j,i] = \bigl(H^{-2}J^\top\bigr)[j,i] = \frac{1}{h_j^2}J[i,j] = \frac{1}{h_j^2}\partDeriv{r_i}{u_j}{} \]
				式(1)の$\tilde{A}_k(\bm{u})$に関する項に着目すると、以下のように計算される。
				但し以下の式で$\delta_{i,j}$はクロネッカーのデルタを表す。
				\begin{align*}
					\phantom{=} &\sum_{i=1}^3 \sum_{j=1}^3 \left(\partDerivLong{\frac{1}{h_k}\partDeriv{r_i}{u_k}{}\tilde{A}_k(\bm{u})}{u_j}{}\right)\frac{1}{h_j^2}\partDeriv{r_i}{u_j}{} = \sum_{j=1}^3 \frac{1}{h_j^2} \sum_{i=1}^3 \left(\partDerivLong{\frac{1}{h_k}\partDeriv{r_i}{u_k}{}\tilde{A}_k(\bm{u})}{u_j}{}\right)\partDeriv{r_i}{u_j}{} \\
					= &\sum_{j=1}^3 \frac{1}{h_j^2} \sum_{i=1}^3 \left[ \left(\partDerivLong{\frac{1}{h_k}}{u_j}{}\right)\tilde{A}_k(\bm{u})\partDeriv{r_i}{u_k}{}\partDeriv{r_i}{u_j}{} + \frac{1}{h_k}\partDeriv{\tilde{A}_k(\bm{u})}{u_j}{}\partDeriv{r_i}{u_k}{}\partDeriv{r_i}{u_j}{} + \frac{\tilde{A}_k(\bm{u})}{h_k}\partDerivIIHetero{r_i}{u_j}{u_k}\partDeriv{r_i}{u_j}{} \right] \\
					= &\sum_{j=1}^3 \frac{1}{h_j^2} \left[ \left(\partDerivLong{\frac{1}{h_k}}{u_j}{}\right)\tilde{A}_k(\bm{u})\partDeriv{\bm{r}}{u_k}{}\cdot\partDeriv{\bm{r}}{u_j}{} + \frac{1}{h_k}\partDeriv{\tilde{A}_k(\bm{u})}{u_j}{}\partDeriv{\bm{r}}{u_k}{}\cdot\partDeriv{\bm{r}}{u_j}{} + \frac{\tilde{A}_k(\bm{u})}{h_k}\partDerivIIHetero{\bm{r}}{u_j}{u_k}\cdot\partDeriv{\bm{r}}{u_j}{} \right] \\
					= &\sum_{j=1}^3 \frac{1}{h_j^2} \left[ \left(\partDerivLong{\frac{1}{h_k}}{u_j}{}\right)\tilde{A}_k(\bm{u})h_k^2\delta_{k,j} + \frac{1}{h_k}\partDeriv{\tilde{A}_k(\bm{u})}{u_j}{}h_k^2\delta_{k,j} + \frac{\tilde{A}_k(\bm{u})}{h_k}\partDerivIIHetero{\bm{r}}{u_j}{u_k}\cdot\partDeriv{\bm{r}}{u_j}{} \right] \\
					= &\frac{1}{h_k^2} \left[ \left(\partDerivLong{\frac{1}{h_k}}{u_k}{}\right)\tilde{A}_k(\bm{u})h_k^2 + h_k\partDeriv{\tilde{A}_k(\bm{u})}{u_k}{} \right] + \frac{\tilde{A}_k(\bm{u})}{h_k} \sum_{j=1}^3 \frac{1}{h_j^2} \partDerivIIHetero{\bm{r}}{u_j}{u_k}\cdot\partDeriv{\bm{r}}{u_j}{} \\
					= &\frac{1}{h_k}\partDeriv{\tilde{A}_k(\bm{u})}{u_k}{} + \frac{\tilde{A}_k(\bm{u})}{h_k} \sum_{j=1}^3 \underbrace{\frac{1}{h_j^2}\partDerivIIHetero{\bm{r}}{u_j}{u_k}\cdot\partDeriv{\bm{r}}{u_j}{}}_{(2)} \quad \left(\because \text{仮定より}\partDerivLong{\frac{1}{h_k}}{u_j}{} = 0\right) \tag{3}
				\end{align*}
				式(2)を評価する。
				$j=k$のとき次式より0である。
				\[ \partDeriv{\bm{r}}{u_k}{2}\cdot\partDeriv{\bm{r}}{u_k}{} = \frac{1}{2}\partDerivLong{\left(\partDeriv{\bm{r}}{u_k}{}\cdot\partDeriv{\bm{r}}{u_k}{}\right)}{u_k}{} = \frac{1}{2}\partDeriv{h_k^2}{u_k}{} = 0 \]
				$j\neq k$のとき次式を得る。
				\begin{align*}
					\phantom{=} &\frac{1}{h_j^2}\partDerivIIHetero{\bm{r}}{u_j}{u_k}\cdot\partDeriv{\bm{r}}{u_j}{} = \frac{1}{h_j^2}\partDerivIIHetero{\bm{r}}{u_k}{u_j}\cdot\partDeriv{\bm{r}}{u_j}{} \\
					&= \frac{1}{h_j^2}\times\frac{1}{2} \partDerivLong{\left(\partDeriv{\bm{r}}{u_j}{}\cdot\partDeriv{\bm{r}}{u_j}{}\right)}{u_k}{} = \frac{1}{h_j^2}\times\frac{1}{2}\partDeriv{h_j^2}{u_k}{} = \frac{1}{h_j}\partDeriv{h_j}{u_k}{}
				\end{align*}
				これらを式(3)に適用して次式を得る。
				但し以下で$[i+j]$は$i+j$を3で割った余りを表す。
				\[ (3) = \frac{1}{h_k}\partDeriv{\tilde{A}_k(\bm{u})}{u_k}{} + \frac{\tilde{A}_k(\bm{u})}{h_k}\left[ \frac{1}{h_{[k+1]}}\partDeriv{h_{[k+1]}}{u_k}{} + \frac{1}{h_{[k+2]}}\partDeriv{h_{[k+2]}}{u_k}{} \right] = \frac{1}{h_1 h_2 h_3}\partDeriv{h_{[k+1]}h_{[k+2]}\tilde{A}_k(\bm{u})}{u_k}{} \]
			\end{proof}
		\section{回転}
			\begin{shadebox}
				既出の記号の定義は引き継ぐ。
				さらに$T$は$\textrm{C}^2$級であるとし、次式が成り立つとする(例: 円柱座標系, 球座標系)。
				\[ \partDerivIIHetero{\bm{r}}{u_j}{u_i}\cdot\partDeriv{\bm{r}}{u_k}{} = 0 \quad \bigl(\{i,j,k\} = \{1,2,3\}\bigr) \tag{1} \]
				3次元Euclid空間上の$\mathrm{C}^1$級のベクトル場$\bm{A}(\bm{r}) = \bm{i}_1 A_1(\bm{r}) + \bm{i}_2 A_2(\bm{r}) + \bm{i}_3 A_3(\bm{r})$を$u_1,u_2,u_3$直交座標系で表したものを$\bm{j}_1 \tilde{A}_1(\bm{u}) + \bm{j}_2 \tilde{A}_2(\bm{u}) + \bm{j}_3 \tilde{A}_3(\bm{u})$とするとき、次式が成り立つ。
				\begin{align*}
					\phantom{=} &\nabla_\bm{r}\times\bm{A}(\bm{r}) \\
					= &\frac{\bm{j}_1}{h_2 h_3}\left[ \partDeriv{h_3\tilde{A}_3}{u_2}{} - \partDeriv{h_2\tilde{A}_2}{u_3}{} \right] + \frac{\bm{j}_2}{h_3 h_1}\left[ \partDeriv{h_1\tilde{A}_1}{u_3}{} - \partDeriv{h_3\tilde{A}_3}{u_1}{} \right] + \frac{\bm{j}_3}{h_1 h_2}\left[ \partDeriv{h_2\tilde{A}_2}{u_1}{} - \partDeriv{h_1\tilde{A}_1}{u_2}{} \right]　\\
					= &\sum_{i=1}^3 \frac{\bm{j}_i}{h_{[i+1]}h_{[i+2]}}\left[ \partDeriv{h_{[i+2]}\tilde{A}_{[i+2]}}{u_{[i+1]}}{} - \partDeriv{h_{[i+1]}\tilde{A}_{[i+1]}}{u_{[i+2]}}{} \right]
				\end{align*}
				ここに$[i+j]$は$i+j$を3で割った余りを表す。
			\end{shadebox}
			\begin{proof}
				\begin{align*}
					\phantom{=} &\nabla_\bm{r}\times\bm{A}(\bm{r}) = \nabla_\bm{r} \times \sum_{i=1}^3 A_i(\bm{r})\bm{i}_i = \sum_{i=1}^3 \nabla_\bm{r} \times A_i(\bm{r})\bm{i}_i = \sum_{i=1}^3 (\nabla_\bm{r}A_i(\bm{r}))\times\bm{i}_i\quad (\because \text{\ref{ベクトルのスカラー関数倍の回転}}) \\
					= &\sum_{i=1}^3 \sum_{j=1}^3 \frac{\bm{j}_j}{h_j}\partDeriv{A_i(\bm{r})}{u_j}{}\times\bm{i}_i = \sum_{i=1}^3 \sum_{j=1}^3 \frac{\bm{j}_j}{h_j}\left(\partDerivLong{\sum_{k=1}^3 \frac{\tilde{A}_k(\bm{u})}{h_k}\partDeriv{r_i}{u_k}{}}{u_j}{}\right) \times \sum_{l=1}^3 \frac{\bm{j}_l}{h_l}\partDeriv{r_i}{u_l}{} \quad (\because \text{\ref{勾配}, \ref{ベクトルの成分の変換}, \ref{単位ベクトルの変換}})
				\end{align*}
				以下では表記の簡便さのため、$\tilde{A}_k(\bm{u})$を単に$\tilde{A}_k$と書く。
				上式の第$m\in\{1,2,3\}$成分を評価する。
				第$m$成分に寄与するのは$(j,l)=\bigl([m+1],[m+2]\bigr),\bigl([m+2],[m+1]\bigr)$のときであり、計算すると次式を得る。
				\begin{align*}
					\phantom{=} &\sum_{i=1}^3 \left[ \frac{1}{h_{[m+1]}}\left(\partDerivLong{\sum_{k=1}^3 \frac{\tilde{A}_k(\bm{u})}{h_k}\partDeriv{r_i}{u_k}{}}{u_{[m+1]}}{}\right)\frac{1}{h_{[m+2]}}\partDeriv{r_i}{u_{[m+2]}}{} \right. \\
					&\left. - \frac{1}{h_{[m+2]}}\left(\partDerivLong{\sum_{k=1}^3 \frac{\tilde{A}_k(\bm{u})}{h_k}\partDeriv{r_i}{u_k}{}}{u_{[m+2]}}{}\right)\frac{1}{h_{[m+1]}}\partDeriv{r_i}{u_{[m+1]}}{} \right] \\
					&= \frac{1}{h_{[m+1]}h_{[m+2]}}\sum_{k=1}^3\left[ \sum_{i=1}^3\left(\partDerivLong{\frac{\tilde{A}_k}{h_k}\partDeriv{r_i}{u_k}{}}{u_{[m+1]}}{}\right)\partDeriv{r_i}{u_{[m+2]}}{} - \sum_{i=1}^3\left(\partDerivLong{\frac{\tilde{A}_k}{h_k}\partDeriv{r_i}{u_k}{}}{u_{[m+2]}}{}\right)\partDeriv{r_i}{u_{[m+1]}}{} \right] \tag{2}
				\end{align*}
				この式の[ ]内第1項を評価すると次式を得る。
				\begin{align*}
					\phantom{=} &\sum_{k=1}^3\sum_{i=1}^3\left(\partDerivLong{\frac{\tilde{A}_k}{h_k}\partDeriv{r_i}{u_k}{}}{u_{[m+1]}}{}\right)\partDeriv{r_i}{u_{[m+2]}}{} \\
					= &\sum_{k=1}^3 \left[ \left(\partDerivLong{\frac{\tilde{A}_k}{h_k}}{u_{[m+1]}}{}\right)\sum_{i=1}^3\partDeriv{r_i}{u_k}{}\partDeriv{r_i}{u_{[m+2]}}{} + \frac{\tilde{A}_k}{h_k}\sum_{i=1}^3\partDerivIIHetero{r_i}{u_{[m+1]}}{u_k}\partDeriv{r_i}{u_{[m+2]}}{} \right] \\
					= &\sum_{k=1}^3 \left[ \left(\partDerivLong{\frac{\tilde{A}_k}{h_k}}{u_{[m+1]}}{}\right)\partDeriv{\bm{r}}{u_k}{}\cdot\partDeriv{\bm{r}}{u_{[m+2]}}{} + \frac{\tilde{A}_k}{h_k}\partDerivIIHetero{\bm{r}}{u_{[m+1]}}{u_k}\cdot\partDeriv{\bm{r}}{u_{[m+2]}}{} \right] \\
					= &\left(\partDerivLong{\frac{\tilde{A}_{[m+2]}}{h_{[m+2]}}}{u_{[m+1]}}{}\right)h_{[m+2]}^2 + \sum_{k\in\{[m+1],[m+2]\}} \frac{\tilde{A}_k}{h_k}\partDerivIIHetero{\bm{r}}{u_{[m+1]}}{u_k}\cdot\partDeriv{\bm{r}}{u_{[m+2]}}{}
				\end{align*}
				最後の等号成立には式(1)を用いた。
				式(2)の[ ]内第2項も同様に評価すると、式(2)の$1/(h_{[m+1]}h_{[m+2]})$の右側の部分は次式になる。
				\begin{align*}
					\phantom{=} &\left(\partDerivLong{\frac{\tilde{A}_{[m+2]}}{h_{[m+2]}}}{u_{[m+1]}}{}\right)h_{[m+2]}^2 - \left(\partDerivLong{\frac{\tilde{A}_{[m+1]}}{h_{[m+1]}}}{u_{[m+2]}}{}\right)h_{[m+1]}^2 \\
					\phantom{=}	& + \sum_{k\in\{[m+1],[m+2]\}}\frac{\tilde{A}_k}{h_k}\left(\partDerivIIHetero{\bm{r}}{u_{[m+1]}}{u_k}\cdot\partDeriv{\bm{r}}{u_{[m+2]}}{} - \partDerivIIHetero{\bm{r}}{u_{[m+2]}}{u_k}\cdot\partDeriv{\bm{r}}{u_{[m+1]}}{}\right) \tag{3}
				\end{align*}
				この式の第3項の()内を評価する。
				$T$は$\mathrm{C}^2$級であると仮定しているから、微分順序の入れ替えが可能であることに留意する。
				$k=[m+1]$のときは次式になる。
				\begin{align*}
					\phantom{=} &\partDeriv{\bm{r}}{u_{[m+1]}}{2}\cdot\partDeriv{\bm{r}}{u_{[m+2]}}{} - \partDerivIIHetero{\bm{r}}{u_{[m+2]}}{u_{[m+1]}}\cdot\partDeriv{\bm{r}}{u_{[m+1]}}{} \\
					= &\partDerivLong{\left(\partDeriv{\bm{r}}{u_{[m+1]}}{}\cdot\partDeriv{\bm{r}}{u_{[m+2]}}{}\right)}{u_{[m+1]}}{} - \partDeriv{\bm{r}}{u_{[m+1]}}{}\cdot\partDerivIIHetero{\bm{r}}{u_{[m+1]}}{u_{[m+2]}}{} - \partDerivIIHetero{\bm{r}}{u_{[m+2]}}{u_{[m+1]}}\cdot\partDeriv{\bm{r}}{u_{[m+1]}}{} \\
					= &-\partDeriv{\bm{r}}{u_{[m+1]}}{}\cdot\partDerivIIHetero{\bm{r}}{u_{[m+2]}}{u_{[m+1]}}{} - \partDerivIIHetero{\bm{r}}{u_{[m+2]}}{u_{[m+1]}}\cdot\partDeriv{\bm{r}}{u_{[m+1]}}{} \\
					= &-\partDerivLong{\left(\partDeriv{\bm{r}}{u_{[m+1]}}{}\cdot\partDeriv{\bm{r}}{u_{[m+1]}}{}\right)}{u_{[m+2]}}{} = -\partDeriv{h_{[m+1]}^2}{u_{[m+2]}}{}
				\end{align*}
				$k=[m+2]$のときは次式になる。
				\begin{align*}
					\phantom{=} &\partDerivIIHetero{\bm{r}}{u_{[m+1]}}{u_{[m+2]}}\cdot\partDeriv{\bm{r}}{u_{[m+2]}}{} - \partDeriv{\bm{r}}{u_{[m+2]}}{2}\cdot\partDeriv{\bm{r}}{u_{[m+1]}}{} \\
					= &\partDerivIIHetero{\bm{r}}{u_{[m+1]}}{u_{[m+2]}}\cdot\partDeriv{\bm{r}}{u_{[m+2]}}{} - \partDerivLong{\left(\partDeriv{\bm{r}}{u_{[m+2]}}{}\cdot\partDeriv{\bm{r}}{u_{[m+1]}}{}\right)}{u_{[m+2]}}{} + \partDeriv{\bm{r}}{u_{[m+2]}}{}\cdot\partDerivIIHetero{\bm{r}}{u_{[m+2]}}{u_{[m+1]}} \\
					= &\partDerivIIHetero{\bm{r}}{u_{[m+1]}}{u_{[m+2]}}\cdot\partDeriv{\bm{r}}{u_{[m+2]}}{} + \partDeriv{\bm{r}}{u_{[m+2]}}{}\cdot\partDerivIIHetero{\bm{r}}{u_{[m+1]}}{u_{[m+2]}} \\
					= &\partDerivLong{\left(\partDeriv{\bm{r}}{u_{[m+2]}}{}\cdot\partDeriv{\bm{r}}{u_{[m+2]}}{}\right)}{u_{[m+1]}}{} = \partDeriv{h_{[m+2]}^2}{u_{[m+1]}}{}
				\end{align*}
				これらの式を(3)に適用すると次式を得る。
				\begin{align*}
					(3) &= \left(\partDerivLong{\frac{\tilde{A}_{[m+2]}}{h_{[m+2]}}}{u_{[m+1]}}{}\right)h_{[m+2]}^2 + \frac{\tilde{A}_{[m+2]}}{h_{[m+2]}}\partDeriv{h_{[m+2]}^2}{u_{[m+1]}}{} - \left(\partDerivLong{\frac{\tilde{A}_{[m+1]}}{h_{[m+1]}}}{u_{[m+2]}}{}\right)h_{[m+1]}^2 - \frac{\tilde{A}_{[m+1]}}{h_{[m+1]}}\partDeriv{h_{[m+1]}^2}{u_{[m+2]}}{} \\
					&= \partDerivLong{\frac{\tilde{A}_{[m+2]}}{h_{[m+2]}}h_{[m+2]}^2}{u_{[m+1]}}{} - \partDerivLong{\frac{\tilde{A}_{[m+1]}}{h_{[m+1]}}h_{[m+1]}^2}{u_{[m+2]}}{} \\
					&= \partDeriv{h_{[m+2]}\tilde{A}_{[m+2]}}{u_{[m+1]}}{} - \partDeriv{h_{[m+1]}\tilde{A}_{[m+1]}}{u_{[m+2]}}{}
				\end{align*}
				これを式(2)に適用して定理の主張を得る。
			\end{proof}
		\section{Laplacian}
			\begin{shadebox}
				既出の記号の定義は引き継ぐ。
				さらに$T$は$\textrm{C}^2$級であるとし、$\frac{\partial h_i}{\partial u_i} = 0\;(i=1,2,3)$が成り立つとする(例: 円柱座標系, 球座標系)。
				$\phi:\realNumbers^3\to\complexNumbers$は$\realNumbers^3$上で$\mathrm{C}^2$級であるとする。
				$\tilde{\phi}(\bm{u}) \coloneqq \phi(T(\bm{u}))$とする。
				$\nabla^2_\bm{r}\phi(\bm{r})$は$\bm{u}$を用いて次式で表せる。
				\[ \nabla^2_\bm{r}\phi(\bm{r}) = \frac{1}{h_1 h_2 h_3}\left[ \partDerivLong{\left(\frac{h_2 h_3}{h_1}\partDeriv{\tilde{\phi}(\bm{u})}{u_1}{}\right)}{u_1}{} + \partDerivLong{\left(\frac{h_3 h_1}{h_2}\partDeriv{\tilde{\phi}(\bm{u})}{u_2}{}\right)}{u_2}{} + \partDerivLong{\left(\frac{h_1 h_2}{h_3}\partDeriv{\tilde{\phi}(\bm{u})}{u_3}{}\right)}{u_3}{} \right] \]
			\end{shadebox}
			\begin{proof}
				\ref{勾配},\ref{発散}を用いて容易に示せる。
			\end{proof}
	\chapter{円柱座標系}
		\section{計量係数}
			\begin{shadebox}
				\[h_r = 1, \hph = r, h_z = 1\]
			\end{shadebox}
			\begin{proof}
				\quad\par
				位置ベクトルは
				\[\vr = \dthv{r\cosph}{r\sinph}{z}\]
				よって
				\begin{align*}
					h_r &\coloneqq \norm{\partDeriv{\vr}{r}{}} = \norm{\dthv{\cosph}{\sinph}{0}} = 1\\
					\hph &\coloneqq \norm{\partDeriv{\vr}{\varphi}{}} = \norm{\dthv{-r\cosph}{r\sinph}{0}} = r\\
					h_z &\coloneqq \norm{\partDeriv{\vr}{z}{}} = \norm{\dthv{0}{0}{1}} = 1
				\end{align*}
			\end{proof}
		\section{単位ベクトルに関する3次元Euclid空間との関係}
			\begin{shadebox}
				\begin{align*}
					\ir &= \ix\cosph + \iy\sinph\\
					\iph &= -\ix\sinph + \iy\cosph
				\end{align*}
				\begin{align*}
					\ix &= \ir\cosph - \iph\sinph\\
					\iy &= \ir\sinph + \iph\cosph
				\end{align*}
			\end{shadebox}
			\begin{proof}
				\quad\par
				位置ベクトルは
				\[\vr = \dthv{r\cosph}{r\sinph}{z}\]
				よって
				\begin{align}
					\ir &\coloneqq \frac{1}{h_r}\partDeriv{\vr}{r}{} = \dthv{\cosph}{\sinph}{0} = \ix\cosph+\iy\sinph \\
					\iph &\coloneqq \frac{1}{\hph}\partDeriv{\vr}{\varphi}{} = \frac{1}{r}\dthv{-r\sinph}{r\cosph}{0} = -\ix\sinph+\iy\cosph
				\end{align}
				これを逆に解いてみる。\\
				$(1)\times\cosph-(2)\times\sinph$より$\ix = \ir\cosph - \iph\sinph$\\
				$(1)\times\sinph+(2)\times\cosph$より$\iy = \ir\sinph + \iph\cosph$\\
				を得る。
			\end{proof}
		\section{微分演算に関する直交座標系との関係}
			\begin{shadebox}
				\begin{align*}
					\partDeriv{}{x}{} &= \cosph\partDeriv{}{r}{} - \frac{\sinph}{r}\partDeriv{}{\varphi}{}\\
					\partDeriv{}{y}{} &= \sinph\partDeriv{}{r}{} + \frac{\cosph}{r}\partDeriv{}{\varphi}{}
				\end{align*}
			\end{shadebox}
			\begin{proof}
				\quad\par
				スカラー関数$\phi(x,y,z)$を微分することを考える。
				\def \phr {\partDeriv{\phi}{r}{}}
				\def \php {\partDeriv{\phi}{\varphi}{}}
				\def \phz {\partDeriv{\phi}{z}{}}
				$d\phi \coloneqq \phr dr + \php d\varphi + \phz dz$であるから
				\begin{align*}
					\partDeriv{\phi}{x}{} &= \phr\partDeriv{r}{x}{} + \php\partDeriv{\varphi}{x}{} + \phz\partDeriv{z}{x}{}\\
					&= \phr\ppx{\sqrt{x^2+y^2}} + \php\ppx{\varphi}
				\end{align*}
				第1項に関して
				\[\ppx{\sqrt{x^2+y^2}} = \frac{x}{\sqrt{x^2+y^2}} = \frac{r\cosph}{r} = \cosph\]
				第2項に関して
				\[\ppx{\varphi} = \ppx\atanEx{x}{y} = \frac{-y}{x^2+y^2} = \frac{-r\sin{\phi}}{r^2} = -\frac{\sin{\phi}}{r}\]
				以上より
				\[\partDeriv{}{x}{} = \cosph\partDeriv{}{r}{} - \frac{\sinph}{r}\partDeriv{}{\varphi}{}\]
				$y$に関しても全く同様にして導出できる。
			\end{proof}
	\chapter{球座標系}
		\section{計量係数}
			\begin{shadebox}
				\[h_r = 1, \hth = r, \hph = r\sinth\]
			\end{shadebox}
			\begin{proof}
				\quad\par
				位置ベクトルは
				\[\vr = \dthv{r\sinth\cosph}{r\sinth\sinph}{r\costh}\]
				よって
				\begin{align*}
					h_r &\coloneqq \norm{\partDeriv{\vr}{r}{}} = \norm{\dthv{\sinth\cosph}{\sinth\sinph}{\costh}} = 1\\
					\hth &\coloneqq \norm{\partDeriv{\vr}{\theta}{}} = \norm{\dthv{r\costh\cosph}{r\costh\sinph}{-r\sinth}} = r\\
					\hph &\coloneqq \norm{\partDeriv{\vr}{\varphi}{}} = \norm{\dthv{-r\sinth\sinph}{r\sinth\cosph}{0}} = r\sinth
				\end{align*}
			\end{proof}
		\section{単位ベクトルに関する3次元Euclid空間との関係}
			\begin{shadebox}
				\begin{align*}
					\ir &= \ix\sinth\cosph + \iy\sinth\sinph + \iz\costh\\
					\ith &= \ix\costh\cosph + \iy\costh\sinph - \iz\sinth\\
					\iph &= -\ix\sinph + \iy\cosph
				\end{align*}
				\begin{align*}
					\ix &= \ir\sinth\cosph + \ith\costh\cosph -\iph\sinph\\
					\iy &= \ir\sinth\sinph + \ith\costh\sinph + \iph\cosph\\
					\iz &= \ir\costh - \ith\sinth
				\end{align*}
			\end{shadebox}
			\noindent(手間が大きい証明)
			\begin{proof}
				\quad\par
				位置ベクトルは
				\[\vr = \dthv{r\sinth\cosph}{r\sinth\sinph}{r\costh}\]
				よって
				\begin{align}
					\ir &\coloneqq \frac{1}{h_r}\partDeriv{\vr}{r}{} = \dthv{\sinth\cosph}{\sinth\sinph}{\costh} = \ix\sinth\cosph + \iy\sinth\sinph + \iz\costh \\
					\ith &\coloneqq \frac{1}{\hth}\partDeriv{\vr}{\theta}{} = \frac{1}{r}\dthv{r\costh\cosph}{r\costh\sinph}{-r\sinth} = \ix\costh\cosph + \iy\costh\sinph - \iz\sinth \\
					\iph &\coloneqq \frac{1}{\hph}\partDeriv{\vr}{\varphi}{} = \frac{1}{r\sinth}\dthv{-r\sinth\sinph}{r\sinth\cosph}{0} = -\ix\sinph + \iy\cosph
				\end{align}
				これを逆に解いてみる。\\
				$(1)\times\costh-(2)\times\sinth$より$\iz = \ir\costh - \ith\sinth$\\
				$(1)\times\sinth+(2)\times\costh$より
				\begin{equation}
					\ir\sinth+\ith\costh = \ix\cosph+\iy\sinph
				\end{equation}
				$(4)\times\cosph-(3)\times\sinph$より$\ix = \ir\sinth\cosph + \ith\costh\cosph -\iph\sinph$\\
				$(3)\times\cosph+(4)\times\sinph$より$\iy = \ir\sinth\sinph + \ith\costh\sinph + \iph\cosph$
			\end{proof}
			\noindent(手間が少ない証明)
			\begin{proof}
				\quad\par
				$\ir,\ith,\iph$の導出までは上の証明と同じ。
				$\ir,\ith,\iph$が直交系を成すことを利用し、$\ix = (\ix\cdot\ir)\ir + (\ix\cdot\ith)\ith + (\ix\cdot\iph)\iph,\quad \iy = (\iy\cdot\ir)\ir + (\iy\cdot\ith)\ith + (\iy\cdot\iph)\iph,\quad \iz = (\iz\cdot\ir)\ir + (\iz\cdot\ith)\ith + (\iz\cdot\iph)\iph$として導出できる。
			\end{proof}
		\section{Laplacian}
			\begin{shadebox}
				球座標系に於けるLaplace演算子は次式である。
				\[ \Delta = \frac{1}{r^2}\frac{\partial}{\partial r}r^2\frac{\partial}{\partial r} + \frac{1}{r^2}\frac{\partial^2}{\partial\theta^2} + \frac{\cos\theta}{r^2\sin\theta}\frac{\partial}{\partial\theta} + \frac{1}{r^2\sin^2\theta}\frac{\partial^2}{\partial\varphi^2} \]
			\end{shadebox}
			\begin{proof}
				\quad\par
				直角座標を引数に取る関数$\phi:\realNumbers^3\to\complexNumbers$は$\mathrm{C}^2$級であるとする。
				球座標から直角座標への変換関数を$f:[r,\theta,\varphi]^\top\in\realNumbers^3\mapsto[r\sin\theta\cos\varphi,r\sin\theta\sin\varphi,r\cos\theta]^\top\in\realNumbers^3$とする。
				$\hat{\phi}(\bm{u}) := \phi(f(\bm{u}))$とする。
				$\bm{r} := [x,y,z]^\top\in\realNumbers^3$とする。
				\begin{align*}
					\Delta_\bm{r}\phi(\bm{r}) &= \frac{\partial^2 \phi}{\partial x^2} + \frac{\partial^2 \phi}{\partial y^2} + \frac{\partial^2 \phi}{\partial z^2} \\
					\begin{split}
						&= \frac{\partial }{\partial x}\left(\frac{\partial \hat{\phi}}{\partial r}\frac{\partial r}{\partial x} + \frac{\partial \hat{\phi}}{\partial \theta}\frac{\partial \theta}{\partial x} + \frac{\partial \hat{\phi}}{\partial \varphi}\frac{\partial \varphi}{\partial x}\right) + \frac{\partial }{\partial y}\left(\frac{\partial \hat{\phi}}{\partial r}\frac{\partial r}{\partial y} + \frac{\partial \hat{\phi}}{\partial \theta}\frac{\partial \theta}{\partial y} + \frac{\partial \hat{\phi}}{\partial \varphi}\frac{\partial \varphi}{\partial y}\right) \\
						&\phantom{=} + \frac{\partial }{\partial z}\left(\frac{\partial \hat{\phi}}{\partial r}\frac{\partial r}{\partial z} + \frac{\partial \hat{\phi}}{\partial \theta}\frac{\partial \theta}{\partial z} + \frac{\partial \hat{\phi}}{\partial \varphi}\frac{\partial \varphi}{\partial z}\right)
					\end{split} \tag{1}
				\end{align*}
				右辺第1項を展開すると次式を得る。
				\begin{align*}
					&\phantom{=} \frac{\partial^2\phi}{\partial r^2}\left(\frac{\partial r}{\partial x}\right)^2 + \frac{\partial \hat{\phi}}{\partial r}\frac{\partial^2 r}{\partial x^2} + \frac{\partial^2\phi}{\partial \theta^2}\left(\frac{\partial \theta}{\partial x}\right)^2 + \frac{\partial \hat{\phi}}{\partial \theta}\frac{\partial^2 \theta}{\partial x^2} + \frac{\partial^2\phi}{\partial \varphi^2}\left(\frac{\partial \varphi}{\partial x}\right)^2 + \frac{\partial \hat{\phi}}{\partial \varphi}\frac{\partial^2 \varphi}{\partial x^2} \\
					&\phantom{=} + 2\frac{\partial^2 \hat{\phi}}{\partial \theta \partial r}\frac{\partial \theta}{\partial x}\frac{\partial r}{\partial x} + 2\frac{\partial^2 \hat{\phi}}{\partial \varphi \partial \theta}\frac{\partial \varphi}{\partial x}\frac{\partial \theta}{\partial x} + 2\frac{\partial^2 \hat{\phi}}{\partial r \partial \varphi}\frac{\partial r}{\partial x}\frac{\partial \varphi}{\partial x}
				\end{align*}
				式(1)の右辺第2,3項はそれぞれ上式の$x$を$y,z$で置き換えたものになる。
				以上より式(1)は次式になる。
				\begin{align*}
					(1) &= \frac{\partial^2 \hat{\phi}}{\partial r^2}\left[\left(\frac{\partial r}{\partial x}\right)^2 + \left(\frac{\partial r}{\partial y}\right)^2 + \left(\frac{\partial r}{\partial z}\right)^2\right] + \frac{\partial \hat{\phi}}{\partial r}\left(\frac{\partial^2 r}{\partial x^2} + \frac{\partial^2 r}{\partial y^2} + \frac{\partial^2 r}{\partial z^2}\right) \tag{2} \\
					&\phantom{=} + \frac{\partial^2 \hat{\phi}}{\partial \theta^2}\left[\left(\frac{\partial \theta}{\partial x}\right)^2 + \left(\frac{\partial \theta}{\partial y}\right)^2 + \left(\frac{\partial \theta}{\partial z}\right)^2\right] + \frac{\partial \hat{\phi}}{\partial \theta}\left(\frac{\partial^2 \theta}{\partial x^2} + \frac{\partial^2 \theta}{\partial y^2} + \frac{\partial^2 \theta}{\partial z^2}\right) \tag{3} \\
					&\phantom{=} + \frac{\partial^2 \hat{\phi}}{\partial \varphi^2}\left[\left(\frac{\partial \varphi}{\partial x}\right)^2 + \left(\frac{\partial \varphi}{\partial y}\right)^2 + \left(\frac{\partial \varphi}{\partial z}\right)^2\right] + \frac{\partial \hat{\phi}}{\partial \varphi}\left(\frac{\partial^2 \varphi}{\partial x^2} + \frac{\partial^2 \varphi}{\partial y^2} + \frac{\partial^2 \varphi}{\partial z^2}\right) \tag{4} \\
					&\phantom{=} + 2\frac{\partial^2 \hat{\phi}}{\partial \theta \partial r}\left(\frac{\partial \theta}{\partial x}\frac{\partial r}{\partial x} + \frac{\partial \theta}{\partial y}\frac{\partial r}{\partial y} + \frac{\partial \theta}{\partial z}\frac{\partial r}{\partial z}\right) \tag{5} \\
					&\phantom{=} + 2\frac{\partial^2 \hat{\phi}}{\partial \varphi \partial \theta}\left(\frac{\partial \varphi}{\partial x}\frac{\partial \theta}{\partial x} + \frac{\partial \varphi}{\partial y}\frac{\partial \theta}{\partial y} + \frac{\partial \varphi}{\partial z}\frac{\partial \theta}{\partial z}\right) \tag{6} \\
					&\phantom{=} + 2\frac{\partial^2 \hat{\phi}}{\partial r \partial \varphi}\left(\frac{\partial r}{\partial x}\frac{\partial \varphi}{\partial x} + \frac{\partial r}{\partial y}\frac{\partial \varphi}{\partial y} + \frac{\partial r}{\partial z}\frac{\partial \varphi}{\partial z}\right) \tag{7}
				\end{align*}
				まず式(5),(6),(7)が0であることを示す。
				$J$を$f$のJacobi行列とすると、式(5)の括弧の内側は$J^{-1}$の1行目と2行目の内積である。
				球座標系は直交座標系であるため$J$の列は直交系を成す。
				これと\ref{行列の列(行)の直交性と逆行列の行(列)の直交性}より$J^{-1}$の行は直交系を成すから、式(5)の括弧の内側は0である。
				同様に式(6),(7)も0である。
				\par
				式(2),(3),(4)を評価する。
				$r = \sqrt{x^2+y^2+z^2},\;\theta = \acos{z/r} = \acos{z/\sqrt{x^2+y^2+z^2}},\;\phi = \atanEx{x}{y}$
				であるから次式が成り立つ。
				\begin{gather*}
					\frac{\partial r}{\partial x} = \sin\theta\cos\varphi, \quad \frac{\partial r}{\partial y} = \sin\theta\sin\varphi, \quad \frac{\partial r}{\partial z} = \cos\theta \\
					\frac{\partial \theta}{\partial x} = \frac{1}{r}\cos\theta\cos\varphi, \quad \frac{\partial \theta}{\partial y} = \frac{1}{r}\cos\theta\sin\varphi, \quad \frac{\partial \theta}{\partial z} = -\frac{1}{r}\sin\theta \\
					\frac{\partial \varphi}{\partial x} = -\frac{\sin\varphi}{r\sin\theta}, \quad \frac{\partial \varphi}{\partial y} = \frac{\cos\varphi}{r\sin\theta}, \quad \frac{\partial \varphi}{\partial z} = 0 \\
					\frac{\partial^2 r}{\partial x^2} = \frac{1}{r}\left(\cos^2\theta\cos^2\varphi + \sin^2\varphi\right), \quad \frac{\partial^2 r}{\partial y^2} = \frac{1}{r}\left(\cos^2\theta\sin^2\varphi + \cos^2\varphi\right), \quad \frac{\partial^2 r}{\partial z^2} = \frac{1}{r}\sin^2\theta \\
					\frac{\partial^2 \theta}{\partial x^2} = \frac{1}{r^2\sin\theta}\left(-2\sin^2\theta\cos\theta\cos^2\varphi + \cos\theta\sin^2\varphi\right) \\
					\frac{\partial^2 \theta}{\partial y^2} = \frac{1}{r^2\sin\theta}\left(-2\sin^2\theta\cos\theta\sin^2\varphi + \cos\theta\cos^2\varphi\right) \\
					\frac{\partial^2 \theta}{\partial z^2} = \frac{2}{r^2}\sin\theta\cos\theta \\
					\frac{\partial^2 \varphi}{\partial x^2} = \frac{2\sin\varphi\cos\varphi}{r^2\sin^2\theta}, \quad \frac{\partial^2 \varphi}{\partial y^2} = -\frac{2\sin\varphi\cos\varphi}{r^2\sin^2\theta}
				\end{gather*}
				これらを用いて
				\begin{align*}
					(2) &= \frac{\partial^2 \hat{\phi}}{\partial r^2} + \frac{2}{r}\frac{\partial \hat{\phi}}{\partial r} = \frac{1}{r^2}\frac{\partial}{\partial r}r^2\frac{\partial \hat{\phi}}{\partial r} \\
					(3) &= \frac{1}{r^2}\frac{\partial^2 \hat{\phi}}{\partial \theta^2} + \frac{\cos\theta}{r^2\sin\theta}\frac{\partial \hat{\phi}}{\partial \theta} \\
					(4) &= \frac{1}{r^2\sin^2\theta}\frac{\partial^2 \hat{\phi}}{\partial \varphi^2}
				\end{align*}
			\end{proof}
